\section{Gott der mich sieht}
Jedes Jahr wird ein bestimmter Vers aus der Bibel als Jahresvers sogenannte \textbf{Jahreslosung} ausgewählt. Für das Jahr 2024 war es der Vers \biblezit{IMos}{16:13}. Dort heisst es: \textit{Du bist ein Gott, der mich sieht.}
Dieser Vers kommt in der Bibel nur an dieser Stelle direkt so vor. Sonst wird in der ganzen Bibel Gott nie so genannt. Wieso reden wir heute über diesen Namen Gottes, wenn er nur einmal in der Bibel vorkommt?

Dieser Name zeigt ein Wesenszug von Gott, der sich durch die ganze Bibel durchzieht. 

Die ganze Geschichte, könnt ihr in der Bibel in \biblezit{IMos}{16:1-16} nachlesen. Nur kurz:

Sahra die Frau von Abraham bekam keine Kinder, wollte aber unbedingt ein Kind. So nötigte sie Abraham mit ihrer Sklavin Hagar zu schlafen, obwohl Gott ihm Nachwuchs prophezeit hat und prompt wurde diese Sklavin schwanger. Die Sahra wurde eifersüchtig und hat ihre Sklavin Hagar weggejagt. Hagar stammte ursprünglich aus Ägypten und machte sich darum wieder auf den Weg dorthin. Unterwegs an einem Brunnen, erschien ihr der \betonung[b, p]{Engel des Herrn}. Dieser sagte ihr in \begin{bibelbox}{SCHL}{IMos}{16:9}
    Und der Engel des Herrn sprach zu ihr: Kehre wieder zurück zu deiner Herrin und demütige dich unter ihrer Hand!
\end{bibelbox}
Dann erzählte er ihr noch, was mit ihrem Sohn passieren wird, dass dieser eine grosse Nachkommenschaft haben wird und ein wilder Mensch sein wird. Er wurde zum Stammvater des arabischen Volkes. Sie soll ihren Sohn Ismael nennen.

Wie immer wird am Anfang des Jahres in vielen Gemeinden und Kirchen über die Jahreslosung gepredigt. So natürlich auch über diesen Vers in 2024. Es wurde durchwegs positiv über diesen Vers berichtet, vor allem, wie Gott dich in deine Not sieht und dir hilft. Was ja auch nicht falsch ist. 

Ich weiss nicht, wenn ich diese Geschichte lese, finde ich nicht, dass dieser Frau in dem Moment wirklich etwas Gutes von Gott getan wurde. Der Engel des Herrn, sagte ihr: \enquote{demütige dich unter ihrer Hand} also der Hand ihrer Herrin Sahra. Sie war doch jetzt frei? Unterwegs nach Hause? Wieso soll es gut sein wieder als Sklavin zurückzukehren?

Wenn wir ihre Geschichte weiter verfolgen, lesen wir, dass es bei ihrer Herrin nicht besser wurde, sondern schlimmer. 17 Jahre später wurde sie entgültig mit ihrem Sohn weggejagt. Diesemal ist sie und ihr Sohn in der Wüste fast umgekommen, aber wieder erschien ihr der Engel des Herrn und hat ihr und ihrem Sohn das Leben gerettet. An dieser Stelle wäre dieser Ausspruch passender gewesen.

Wir sehen \betonung[b, p]{Gott ist ein Gott, der dich sieht}, Hagar hat es zweimal erlebt. Beide Male war Hagar dem Engel des Herrn gehorsam. So gesehen war sie Gott gegenüber gehorsamer als Sahra.

Liest die Geschichte von Hiob. Im ersten Kapitel lesen wir, dass Gott Hiob sieht in \biblezit{Hi}{1:8}. Ging es Hiob danach gut? Paolo wird euch in seiner Botschaft mitteilen, wie es Hiob ergangen ist.

Ich weiss, Gott sieht mich, Gott achtet auf mich, Gott gebraucht mich. Davon bin ich überzeugt. Er sieht nicht nur wiedergeborene Christen, er sieht alle Menschen und gebraucht alle. Bei der 3. Rückführung der Juden nach Israel ab 1880, waren mehr Atheisten beteiligt als gläubige Menschen. Theodor Herzl \zB, der Vater des Zionismus, war Atheist.

Was ich aber habe, was ein Atheist nicht hat, ist die Gewissheit, dass ich ein ewiges Leben zusammen mit diesem \betonung[b, p]{Gott, der mich sieht}, verbringen werde. Wieso habe ich die Gewissheit? Weil der Engel des Herrn, der im Neuen Testament mit dem Namen Jesus als Mensch auf die Erde kam, mir es versprochen hat. Dafür liess sich Jesus für uns opfern. Genau so, wie vor 6000 Jahren die Prophezeiung, dass aus Ismael ein riesiges Volk werden wird, wahr wurde, dürfen auch wir auf sein Versprechen welches er uns in \biblezit{Joh}{5:24} gab, zu 100 \% vertrauen. Ist das nicht ein Trost in unserem kurzen Dasein hier auf der Erde? Nun woher wissen wir auf welches Wort wir hören sollen? Als Jesus das sagte, gab noch keinen Papst, Pastor oder Priester? Nur durch sein Wort aus der Bibel, die wir hier haben. \betonung[b, p]{Dieses Wort ist der einzige Massstab}.

Israel ist das Volk Gottes und bekommt von ihm eine Sonderbehandlung hier auf Erden. Wir wiedergeborene Christen, also die, die glauben, dass Jesus Christus für unsere Sünden am Kreuz gestorben ist und wieder den Toten auferstanden ist, sind die Gemeinde Jesus, egal ob katholisch, reformiert oder vegan. Auch wenn wir wissen \betonung[b, p]{das wir einem Gott danken, der uns sieht}, wissen wir, dass wir in einer gefallenen Welt leben.\\
Oder was machten die 12 Christen im Kongo vor Gott falsch, dass sie wegen ihres Glaubens geköpft wurden? \betonung[b, p]{Der Gott, der dich sieht}, hat auch dies gesehen. Wir werden die Sonderbehandlung erst im Himmel haben. Bis dahin, denke immer daran, Jesus ist für deine und meine Sünden am Kreuz gestorben, er hat gelitten für dich und mich und er ist Gott der dich und mich sieht. Weil das so ist, achten wir doch darauf, dass wir möglichst ein würdevolles, dankbares Leben leben, um sein Opfer für uns und ihn zu ehren. Darum werden wir heute zu Gendenken an sein Opfer für uns das Abendmahl feiern, wie er es uns aufgetragen hat. Nicht vergessen \betonung[b, p]{Er ist ein Gott, der dich und mich sieht} IMMER UND ÜBERALL.